\subsection{Direct/Indirect Prompt Injection 与 Memory/RAG Poisoning}

slides 39--50 从 direct prompt injection 走向 indirect injection。Direct attack 在用户输入中要求忽略系统指令、泄露 prompt 或执行目标行为；indirect attack 把指令藏在网页、邮件、文档、候选人材料或检索结果里，让 Agent 在“读取数据”时误把数据当命令。攻击面还包括 memory poisoning、knowledge-base poisoning、tool description 与多 Agent 消息。AgentPoison 进一步展示了对 RAG memory/knowledge base 的后门化。

\lecturefigure{slides-images/slide-039.png}{官方 slide 39：Direct Prompt Injection}
\lecturefigure{slides-images/slide-040.png}{官方 slide 40：Bing Chat system prompt leakage}
\lecturefigure{slides-images/slide-041.png}{官方 slide 41：Prompt Injection attack methods}
\lecturefigure{slides-images/slide-042.png}{官方 slide 42：Indirect Prompt Injection 场景}
\lecturefigure{slides-images/slide-044.png}{官方 slide 44：恶意内容进入检索与决策链}
\lecturefigure{slides-images/slide-048.png}{官方 slide 48：Indirect injection 的最终错误决策}
\lecturefigure{slides-images/slide-049.png}{官方 slide 49：Prompt Injection attack surface}
\lecturefigure{slides-images/slide-050.png}{官方 slide 50：AgentPoison 对 RAG memory/knowledge base 的后门攻击}

\begin{knowledgebox}{读图：Data 与 Instruction 必须有不可混淆的 provenance}
Indirect injection 的根因是系统把不可信数据提升为可执行指令。仅在 prompt 中写“不要听网页的话”不构成强边界；更稳妥的做法是给来源加标签，限制外部内容只能进入 data channel，策略层决定哪些字段可影响计划，tool 层再验证动作是否符合用户目标和权限。
\end{knowledgebox}

\begin{warningbox}{Memory poisoning 让一次攻击变成长期攻击}
若恶意内容被摘要后写入长期 memory，后续会话即使没有原始 payload 也可能受影响。写入 memory 前要验证来源、用途和时效，保留原始证据与撤销机制，并对高影响记忆要求用户或独立 policy 批准。
\end{warningbox}

